\SourceLeaf{068}
\[
=\frac{1}{\dfrac{LL}{p}+\dfrac{L'L'}{p'}+\dfrac{L''L''}{p''}+\text{etc.}}\quad(3)
\]
sive \(\frac1P\) erit valor polynomii supra allati pro eo systemate valorum quantitatum \(x,y,z\) etc., per quod aequationibus (1) satisfit.

\Article{6}
In art. praec. eam functionem \(U\) dignoscere docuimus, quae determinationi maxime idoneae incognitae nostrae inservit: videamus iam, quemnam \textit{valorem} incognita hoc modo assequatur. Designetur hic valor per \(K\), qui itaque oritur, si in \(U\) valores observati quantitatum \(v,v',v''\) etc. substituuntur; per eandem substitutionem obtineat functio \(u\) valorem \(k\); denique sit \(\kappa\) valor verus incognitae, qui proin e valoribus veris quantitatum \(v,v',v''\) etc. proditurus esset, si hos vel in \(U\) vel in \(u\) substituere possemus. Hinc itaque erit
\[
\begin{aligned}
k&=\kappa+le+l'e'+l''e''+\text{etc.}\\
K&=\kappa+Le+L'e'+L''l''+\text{etc.}
\end{aligned}
\]
adeoque
\[
K=k(L-l)e+(L'-l')e'+(L''-l'')e''+\text{etc.}
\]
Substituendo in hac aequatione pro \(L-l,\ L'-l',\ L''-l''\) etc. valores ex (2), statuendoque
\[
\left.\begin{aligned}
ae+a'e'+a''e''+\text{etc.}&=\mathfrak A\\
be+b'e'+b''e''+\text{etc.}&=\mathfrak B\\
ce+c'e'+c''e''+\text{etc.}&=\mathfrak C
\end{aligned}\right\}\ (4)
\]
etc., habebimus
\[
K=k+\mathfrak A x^0+\mathfrak B y^0+\mathfrak C z^0\ \text{etc.}\qquad(5)
\]
Valores quantitatum \(\mathfrak A,\mathfrak B,\mathfrak C\) etc. per formulas (4) quidem calculare non possumus, quum errores \(e,e',e''\) etc. maneant incogniti; at sponte manifestum est, illos nihil aliud esse, nisi valores functionum \(X,Y,Z\) etc., qui prodeunt, si pro \(v,v',v''\) etc. valores observati substituuntur. Hoc modo systema aequationum (1), (3), (5) completam problematis nostri solutionem exhibet, quum ea, quae in fine art. 2. de computo quantitatum \(l,l',l''\) etc, valoribus observatis quantita\SourceLeaf{069}%
tum \(v,v',v''\) etc. superstruendo monuimus, manifesto aequali iure ad computum quantitatum \(a,a',a''\) etc. \(b,b',b''\) etc. etc. extendere liceat.

\Article{7}
Loco formulae (3), pondus determinationis maxime plausibilis exprimentis, plures aliae exhiberi possunt, quas evolvere operae pretium erit.

Primo observamus, si aequationes (2) resp. per \(\frac ap,\ \frac{a'}{p'},\ \frac{a''}{p''}\) etc. multiplicentur et addantur. prodire
\[\relax
[aa]x^0+[ab]y^0+[ac]z^0+\text{etc.}=\frac{aL}{p}+\frac{a'L'}{p'}+\frac{a''L''}{p''}+\text{etc.}
\]
Pars ad laevam fit \(=0\), partem ad dextram iuxta analogiam per \([aL]\) denotamus; habemus itaque
\[\relax
[aL]=0,\quad\text{et prorsus simili modo}\quad[bL]=0,\ [cL]=0\ \text{etc.}
\]
Multiplicando porro aequationes (2) deinceps per \(\frac Lp,\ \frac{L'}{p'},\ \frac{L''}{p''}\) etc., et addendo, invenimus
\[
\frac{lL}{p}+\frac{l'L'}{p'}+\frac{l''L''}{p''}+\text{etc.}=\frac{LL}{p}+\frac{L'L'}{p'}+\frac{L''L''}{p''}+\text{etc.}
\]
unde obtinemus expressionem \textit{secundam} pro pondere,
\[
P=\frac{1}{\dfrac{lL}{p}+\dfrac{l'L'}{p'}+\dfrac{l''L''}{p''}+\text{etc.}}
\]
Denique multiplicando aequationes (2) deinceps per \(\frac lp,\ \frac{l'}{p'},\ \frac{l''}{p''}\) etc., et addendo, pervenimus ad expressionem \textit{tertiam} ponderis
\[
P=\frac{1}{[al]x^0+[bl]y^0+[cl]z^0+\text{etc.}+[ll]}
\]
si ad instar reliquarum denotationum statuimus
\[
\frac{ll}{p}+\frac{l'l'}{p'}+\frac{l''l''}{p''}+\text{etc.}=[ll]
\]
Hinc adiumento aequationum (1) facile fit transitus ad \textit{expressionem quartam}, quam ita exhibemus:
\[
\begin{aligned}
\frac1P=[ll]&-[aa]x^0x^0-[bb]y^0y^0-[cc]z^0z^0-\text{etc.}\\
&-2[ab]x^0y^0-2[ac]x^0z^0-2[bc]y^0z^0-\text{etc.}
\end{aligned}
\]
